\scanpage{121}

% manuscript lines / position: 1–5
$\unsure{\mathscr X|\mathcal X}$ sous $\unsure{\gamma|\mathcal G|\Gamma}$, et de la transitivité de $\gamma$ sur \unsure{$I$|$J$}, on voit que cet idéal $\unsure{\mathfrak a|\mathcal A|\mathfrak J}_j$ ne dépend pas de $j$. D'autre part, l'intersection des $\mathfrak a_j$ est évidemment le noyau de $\deleted{\unsure{\mathfrak a}}k\longrightarrow\Gamma(X_J,\unsure{\underline O|\mathcal O}_{X_J})\simeq k^J$. Comme celui-ci est nul, on a $\mathfrak a_j=0$ donc

% manuscript lines / position: 6
\[
 \mathscr X_j=X_j\quad\text{donc}\quad\mathscr X=X_J\qquad\text{O.K.}
\]

% manuscript lines / position: 7–13
Ce résultat vaut sans restriction de caractéristique, mais il concerne la situation standard stricte, \uline{très} restrictive. J'ai envie de vérifier que dès qu'on a trois rev. quadratiques $X_1,X_2,X_3$ de $X$, et une section $\varepsilon$ de $X_1\wedge X_2\wedge X_3$, sans condition étale, on en déduit sur

% manuscript lines / position: 14
\[
 X=\pi^{-1}(\varepsilon)\qquad
 \deleted{\unsure{\pi}}\; X_1\times_X X_2\times_X X_3\xrightarrow{\pi}X_1\wedge X_2\wedge X_3
\]

% manuscript lines / position: 15–17
une structure biquadratique p.r. aux $\unsure{X_i|A_i}$. Ceci doit généraliser la structure biquadratique can. sur un $X_1\times X_2$, par les applications

% manuscript lines / position: 18–20
\diagram{\begin{tikzcd}[ampersand replacement=\&,row sep=small]
 \& X_1\\
 X_1\times X_2 \arrow[ur] \arrow[r] \arrow[dr] \& X_2\\
 \& X_3=X_1\wedge X_2,
\end{tikzcd}}

% manuscript lines / position: 21
comme on va le voir \unsure{tantôt|bientôt|de suite}.

% manuscript lines / position: bottom left
% archive sheet number: 120
